\documentclass[spanish,openany]{report}
% \usepackage{darkmode} \enabledarkmode
\usepackage{wrapfig}
\usepackage{enumitem}
\input{~/git/Clase/template/preamble.tex}
\input{~/git/Clase/template/macros.tex}
\input{~/git/Clase/template/letterfonts.tex}

\date{\today}
% \hbadness=99999
\usepackage{microtype}
% \addbibresource{./bibliografia.bib}
\usepackage{titling}
\usepackage{afterpage}


% Header and footer
% \fancypagestyle{plain}{
%     \fancyhf{} %Clear Everything.
%     \fancyfoot[C]{Página \thepage\ de \pageref{LastPage} }%Page Number
%     \renewcommand{\headrulewidth}{1pt} %0pt for no rule, 2pt thicker etc...
%     \renewcommand{\footrulewidth}{1pt}
%     \fancyfoot[R]{\small{Daniel Carbajales Soto}}
%     \fancyhead[R]{\small{Cinta Transportadora con 3 Velocidades}}
%     \fancyhead[L]{\small{\today}}
% }
% \pagestyle{plain}

\begin{document}

%----------------------------------------------------------------------------------------
%	                            PORTADA
%----------------------------------------------------------------------------------------
% \makeportada{}{Daniel Carbajales Soto - 1° CFGSARI}{}

%----------------------------------------------------------------------------------------
%                                    HEADER 
%----------------------------------------------------------------------------------------
% \header{}{Daniel Carbajales Soto}{1° CFGSARI}
%----------------------------------------------------------------------------------------
%                               TABLE OF CONTENTS
%----------------------------------------------------------------------------------------
% \maketoc
%----------------------------------------------------------------------------------------
%	                            TABLES
%----------------------------------------------------------------------------------------
% this is just a demo in case i need to make a table
%
%
%
% \begin{tabularx}{\linewidth}{|C{2cm}|Y|Y|C{2cm}|}
% \hline
% \makecell{Header 1} & \makecell{Header 2 \\ with linebreak} & \makecell{Header 3} & \makecell{Header 4} \\
% \hline
% Row 1 & This is a long cell text that automatically wraps in the Y column & Another long cell & Short \\
% \hline
% Row 2 & \multirow{2}{*}{Multirow text} & More text & Short \\
% \cline{1-1} \cline{3-4}
% Row 3 & & Extra text & Short \\
% \hline
% \end{tabularx}


%----------------------------------------------------------------------------------------
%	                            DOCUMENT
%----------------------------------------------------------------------------------------

\chapter*{Ejercicios de protección electrica}


\begin{enumerate}
  \item \boldit{Representa los símbolos de los siguientes elementos}

  \begin{quote}
  % ✍️ Write your response here
  \end{quote}

  \item \boldit{Asocia el símbolo con el dispositivo correspondiente y di cómo se denomina}

  \begin{quote}
	  \begin{itemize}[label=]
		  \item 1 $\to$ B: Magnetotérmico
		  \item 2 y 5 $\to$ D: Fusibles
		  \item 3 $\to$ B: Protección contra sobretensiónes
		  \item 4 $\to$ E: DPN
  	  \end{itemize}
  \end{quote}

  \item \boldit{Dibuja el esquema del circuito de protección de instalación monofásica con las siguientes características}

  \begin{quote}
	  \begin{itemize}
		\item{Debe de disponer de tres lineas independientes}
		\item{Protegidas contra sobreintensidades de manera individual}
		\item{Protección diferencial común}
	  \end{itemize}
	  \incfigpdf[0.35]{ejercicio3_esquema_multifilar}  

  \end{quote}

  \item \boldit{Di cuáles de estas afirmaciones son incorrectas en relación con los interruptores diferenciales y explica por qué}

  \begin{quote}
	  \begin{enumerate}[label=\Roman*)]
		  \item{\boldit{incorrecta.}} No existen.
		  \item[III)]{\boldit{incorrecta.}} No protegen contra cortocircuitos y sobrecargas.
	  \end{enumerate}
  \end{quote}
\clearpage
  \item \boldit{Observa el siguiente cuadro de protección correspondiente a una instalación bifásica a 230 V y dibuja el esquema unifilar del circuito}

  \begin{quote}
	  \incfigpdf[0.8]{esquema_unifilar_1}
  \end{quote}

  \item \boldit{En un local industrial, con alimentación monofásica, se dispone de dos líneas para alumbrado y de dos para fuerza (bases de enchufes). Dibuja cuál es el esquema multifilar de protección, utilizando para cada línea un interruptor magnetotérmico y un solo diferencial para alumbrado y otro para fuerza. Además, debe estar previsto la instalación de un interruptor de corte general de tipo magnetotérmico}
  \begin{quote}
  % ✍️ Write your response here

\incfigpdf[0.5]{SEHN_ejercicio-6}
  \end{quote}

  \item \boldit{Nombra los diferentes valores de sensibilidad de los interruptores diferenciales}
	\begin{quote}
		\begin{itemize}
			\item{30mA}
			\item{100mA}
			\item{300mA}
			\item{>300mA}
		\end{itemize}
	\end{quote}
\clearpage
  \item \boldit{En las características de un interruptor diferencial ¿Qué diferencias hay entre la corriente nominal y la sensibilidad? ¿En qué se miden?}

  \begin{quote}
	  La corriente nominal es el amperaje que soporta el dispositivo medido en \textbf{``Amperios''}. La sensibilidad es la corriente de fuga 
	  que puede detectar el diferencial antes de que se dispare y es medida en \textbf{``miliamperio''} o \textbf{``mA''}.
  \end{quote}

  \item \boldit{Dibuja el esquema multifilar de un cuadro de protección para una línea de alimentación trifásica con neutro y conductor de protección a 400 Vca}

  \begin{quote}
	  	
	  \begin{paracol}{2}
		  \incfigpdf[1]{SEHN_ejercicio-9}
		  \switchcolumn
		\begin{itemize}
			\vfill
			  \item{Un circuito trifásico con neutro a 400v}
			\vfill
			  \item{Un circuito trifásico sin neutro a 400v}
			\vfill
			  \item{Tres lineas monofásicas a 230v}
			\vfill

		\end{itemize}
	  \end{paracol}
  \end{quote}

  \item \boldit{Indica como se realiza la conexión de un protector contra sobretensiones y haz un esquema donde se aprecie}

  \begin{quote}
	  Conectamos las fases, el neutro y la tierra en paralelo con la instalación después del automático general o del diferencial.
	  tal que así:\\
	  \incfigpdf[0.5]{SEHN_ejercicio-10}

  \end{quote}
  \clearpage

  \item \boldit{¿En qué consiste la instalación de puesta a tierra?}

  \begin{quote}
  % ✍️ Write your response here
  \end{quote}

  \item \boldit{¿A qué se denomina conductores de protección?}

  \begin{quote}
  % ✍️ Write your response here
  \end{quote}

  \item \boldit{Indica los tipos de protectores contra sobretensiones y cuál es el que se utiliza con más frecuencia en instalaciones de interior}

  \begin{quote}
  % ✍️ Write your response here
  \end{quote}

  \item \boldit{¿Qué parámetros debemos tener en cuenta a la hora de seleccionar un magnetotérmico?}

  \begin{quote}
  % ✍️ Write your response here
  \end{quote}

  \item \boldit{¿Cuándo normalmente se utilizan fusibles de protección en instalaciones interiores?}

  \begin{quote}
  % ✍️ Write your response here
  \end{quote}

  \item \boldit{Indica cuáles son los principales fallos eléctricos que pueden aparecer en una instalación eléctrica}

  \begin{quote}
  % ✍️ Write your response here
  \end{quote}

  \item \boldit{En una línea eléctrica de 400 V de tres fases con neutro, ¿qué tensiones se medirán entre las siguientes opciones?}

  \begin{quote}
  % ✍️ Write your response here
  \end{quote}

  \item \boldit{¿Contra qué anomalías actúan las protecciones que se ponen en las instalaciones eléctricas?}

  \begin{quote}
  % ✍️ Write your response here
  \end{quote}

  \item \boldit{¿Qué es un cortocircuito y por qué es tan peligroso?}

  \begin{quote}
  % ✍️ Write your response here
  \end{quote}

  \item \boldit{¿Cómo se comprueba si un fusible está o no fundido? ¿Cómo se interpreta dicha comprobación?}

  \begin{quote}
  % ✍️ Write your response here
  \end{quote}

  \item \boldit{¿Para qué sirve el botón de prueba de un diferencial? ¿En qué condiciones actúa?}

  \begin{quote}
  % ✍️ Write your response here
  \end{quote}

  \item \boldit{Di en qué se parecen y en qué se diferencian los fusibles y los magnetotérmicos}

  \begin{quote}
  % ✍️ Write your response here
  \end{quote}

  \item \boldit{¿Qué diferencias hay entre una sobrecarga y un cortocircuito?}

  \begin{quote}
  % ✍️ Write your response here
  \end{quote}

  \item \boldit{¿Qué es un contacto indirecto? Represéntalo gráficamente}

  \begin{quote}
  % ✍️ Write your response here
  \end{quote}

  \item \boldit{¿A qué se denomina corriente de fuga o derivación?}

  \begin{quote}
  % ✍️ Write your response here
  \end{quote}

  \item \boldit{Observa la siguiente lista de anomalías que se pueden producir en una instalación eléctrica y di con qué dispositivo o dispositivos se pueden proteger}

  \begin{quote}
  % ✍️ Write your response here
  \end{quote}

  \item \boldit{¿Qué misión cumple el diferencial en el circuito de protección de una instalación eléctrica? ¿El diferencial se dispara ante sobrecorrientes?}

  \begin{quote}
  % ✍️ Write your response here
  \end{quote}

  \item \boldit{¿Qué es un contacto directo?}

  \begin{quote}
  % ✍️ Write your response here
  \end{quote}

  \item \boldit{¿Cuáles son los tipos de sobretensiones que pueden producirse en función del tiempo?}

  \begin{quote}
  % ✍️ Write your response here
  \end{quote}

  \item \boldit{¿Cuál es el objetivo de la puesta a tierra?}

  \begin{quote}
  % ✍️ Write your response here
  \end{quote}

  \item \boldit{¿Por qué es fundamental la instalación de puesta a tierra para el correcto funcionamiento de los diferenciales?}

  \begin{quote}
  % ✍️ Write your response here
  \end{quote}

  \item \boldit{Explica las dos formas de disparo de un interruptor magnetotérmico}

  \begin{quote}
  % ✍️ Write your response here
  \end{quote}
\end{enumerate}





%%%%%%%%%%%%%%%%%%%%%%%%%%%%%%%%%%%%%%%%%%%%%%%%%%%%%%%%%%
%%%%%%%%%%%%%%%%%%%% bibliography %%%%%%%%%%%%%%%%%%%%%%%%
% \nocite{*}                                 %%%%%%%%%%%%%
% \printbibliography                         %%%%%%%%%%%%%
%%%%%%%%%%%%%%%%%%%%%%%%%%%%%%%%%%%%%%%%%%%%%%%%%%%%%%%%%%
%%%%%% Last page in case there is no bibliography %%%%%%%%
% \label{Last Page}                           %%%%%%%%%%%%%
%%%%%%%%%%%%%%%%%%%%%%%%%%%%%%%%%%%%%%%%%%%%%%%%%%%%%%%%%%

\end{document}
